% Options for packages loaded elsewhere
\PassOptionsToPackage{unicode}{hyperref}
\PassOptionsToPackage{hyphens}{url}
\documentclass[
  11pt,
]{article}
\usepackage{xcolor}
\usepackage[left=3cm,right=5cm,top=2cm,bottom=2cm]{geometry}
\usepackage{amsmath,amssymb}
\setcounter{secnumdepth}{5}
\usepackage{iftex}
\ifPDFTeX
  \usepackage[T1]{fontenc}
  \usepackage[utf8]{inputenc}
  \usepackage{textcomp} % provide euro and other symbols
\else % if luatex or xetex
  \usepackage{unicode-math} % this also loads fontspec
  \defaultfontfeatures{Scale=MatchLowercase}
  \defaultfontfeatures[\rmfamily]{Ligatures=TeX,Scale=1}
\fi
\usepackage{lmodern}
\ifPDFTeX\else
  % xetex/luatex font selection
\fi
% Use upquote if available, for straight quotes in verbatim environments
\IfFileExists{upquote.sty}{\usepackage{upquote}}{}
\IfFileExists{microtype.sty}{% use microtype if available
  \usepackage[]{microtype}
  \UseMicrotypeSet[protrusion]{basicmath} % disable protrusion for tt fonts
}{}
\makeatletter
\@ifundefined{KOMAClassName}{% if non-KOMA class
  \IfFileExists{parskip.sty}{%
    \usepackage{parskip}
  }{% else
    \setlength{\parindent}{0pt}
    \setlength{\parskip}{6pt plus 2pt minus 1pt}}
}{% if KOMA class
  \KOMAoptions{parskip=half}}
\makeatother
\usepackage{longtable,booktabs,array}
\usepackage{caption}
\captionsetup[table]{skip=6pt}
\newcounter{none} % for unnumbered tables
\usepackage{calc} % for calculating minipage widths
% Correct order of tables after \paragraph or \subparagraph
\usepackage{etoolbox}
\makeatletter
\patchcmd\longtable{\par}{\if@noskipsec\mbox{}\fi\par}{}{}
\makeatother
% Allow footnotes in longtable head/foot
\IfFileExists{footnotehyper.sty}{\usepackage{footnotehyper}}{\usepackage{footnote}}
\makesavenoteenv{longtable}
\usepackage{graphicx}
\makeatletter
\newsavebox\pandoc@box
\newcommand*\pandocbounded[1]{% scales image to fit in text height/width
  \sbox\pandoc@box{#1}%
  \Gscale@div\@tempa{\textheight}{\dimexpr\ht\pandoc@box+\dp\pandoc@box\relax}%
  \Gscale@div\@tempb{\linewidth}{\wd\pandoc@box}%
  \ifdim\@tempb\p@<\@tempa\p@\let\@tempa\@tempb\fi% select the smaller of both
  \ifdim\@tempa\p@<\p@\scalebox{\@tempa}{\usebox\pandoc@box}%
  \else\usebox{\pandoc@box}%
  \fi%
}
% Set default figure placement to htbp
\def\fps@figure{htbp}
\makeatother
% definitions for citeproc citations
\NewDocumentCommand\citeproctext{}{}
\NewDocumentCommand\citeproc{mm}{%
  \begingroup\def\citeproctext{#2}\cite{#1}\endgroup}
\makeatletter
 % allow citations to break across lines
 \let\@cite@ofmt\@firstofone
 % avoid brackets around text for \cite:
 \def\@biblabel#1{}
 \def\@cite#1#2{{#1\if@tempswa , #2\fi}}
\makeatother
\newlength{\cslhangindent}
\setlength{\cslhangindent}{1.5em}
\newlength{\csllabelwidth}
\setlength{\csllabelwidth}{3em}
\newenvironment{CSLReferences}[2] % #1 hanging-indent, #2 entry-spacing
 {\begin{list}{}{%
  \setlength{\itemindent}{0pt}
  \setlength{\leftmargin}{0pt}
  \setlength{\parsep}{0pt}
  % turn on hanging indent if param 1 is 1
  \ifodd #1
   \setlength{\leftmargin}{\cslhangindent}
   \setlength{\itemindent}{-1\cslhangindent}
  \fi
  % set entry spacing
  \setlength{\itemsep}{#2\baselineskip}}}
 {\end{list}}
\usepackage{calc}
\newcommand{\CSLBlock}[1]{\hfill\break\parbox[t]{\linewidth}{\strut\ignorespaces#1\strut}}
\newcommand{\CSLLeftMargin}[1]{\parbox[t]{\csllabelwidth}{\strut#1\strut}}
\newcommand{\CSLRightInline}[1]{\parbox[t]{\linewidth - \csllabelwidth}{\strut#1\strut}}
\newcommand{\CSLIndent}[1]{\hspace{\cslhangindent}#1}
\setlength{\emergencystretch}{3em} % prevent overfull lines
\providecommand{\tightlist}{%
  \setlength{\itemsep}{0pt}\setlength{\parskip}{0pt}}
% Three-part paragraph scaffolding for manuscript revision.
%
% bullets  -- boiled-down topic sentence + bullet points
%             small, sans-serif, dark gray; paragraph number [N] prepended
% rgt      -- author's rewrite (initially a placeholder)
%             small, italic, dark blue
% orig     -- original Claude-authored draft
%             normal size, light gray
%
% Presentation order per paragraph: bullets -> rgt -> orig
% Strip with: strip-claudecode -i report.Rmd

\usepackage{xcolor}

\newcounter{pgraph}

\newenvironment{bullets}{%
  \stepcounter{pgraph}%
  \begin{quote}\small\sffamily\color[gray]{0.30}%
  {\normalfont\scriptsize\color[gray]{0.58}[\thepgraph]\enspace}}{%
  \end{quote}}

\newenvironment{rgt}{%
  \begin{quote}\small\itshape\color[RGB]{20,60,160}}{%
  \end{quote}}

\newenvironment{orig}{%
  \begin{quote}\normalsize\color[gray]{0.55}}{%
  \end{quote}}

% backward compat alias
\newenvironment{claudecode}{\begin{orig}}{\end{orig}}
\usepackage{amsmath}
\usepackage{amssymb}
\usepackage{lineno}
\linenumbers
\usepackage{setspace}
\setstretch{1.1}
\usepackage{bookmark}
\IfFileExists{xurl.sty}{\usepackage{xurl}}{} % add URL line breaks if available
\urlstyle{same}
% fallback for those not using the hyperref driver hyperxmp:
\makeatletter
\@ifundefined{xmpquote}{\newcommand{\xmpquote}[1]{#1}}{}
\makeatother
\hypersetup{
  pdftitle={Test-procedure choice for the biomarker-treatment interaction in aggregated N-of-1 trials: dichotomization{,} classical RM-ANOVA{,} mixed-effects{,} and generalized estimating equations},
  pdfauthor={pmsimstats team},
  hidelinks,
  pdfcreator={LaTeX via pandoc}}

\title{Test-procedure choice for the biomarker-treatment interaction in aggregated N-of-1 trials: dichotomization, classical RM-ANOVA, mixed-effects, and generalized estimating equations}
\author{pmsimstats team}
\date{2026-07-30}

\begin{document}
\maketitle

{
\setcounter{tocdepth}{2}
\tableofcontents
}
\section{Abstract}\label{abstract}

\textbf{Background.} Inference about the biomarker-treatment
interaction in aggregated N-of-1 trials depends on the analyst's
test procedure: classical repeated-measures ANOVA (RM-ANOVA),
the linear mixed-effects model (LME), or generalized estimating
equations (GEE). A common but rarely examined feature of the
classical analysis is that it requires the continuous biomarker
to be dichotomized. We quantify the consequences of test-procedure
choice for the interaction, separating the cost of dichotomization
from the cost of the test procedure proper.

\textbf{Methods.} We derive the closed-form strict RM-ANOVA
\(F\)-statistic for the biomarker-stratified treatment-by-time
interaction in a balanced split-plot design with categorical
predictors and the sphericity assumption, and relate it to the
Wald \(t\)-test for the corresponding fixed-effects coefficient in
a linear mixed model with continuous-time AR(1) residual
correlation. We then conduct a Monte Carlo simulation study,
designed and reported under the\textsuperscript{\citeproc{ref-morris2019}{1}} ADEMP framework,
that crosses five test procedures (strict RM-ANOVA on the
dichotomized biomarker; a continuous-biomarker within-subject
contrast; the standard pmsimstats linear-mixed analysis with
\texttt{corCAR1}; and GEE with naive and with Mancl-DeRouen small-sample
sandwich variance) against the biomarker-moderation coefficient
\(c_{bm} \in \{0, 0.45\}\) and sample size \(N \in \{25, 35, 70,
100\}\), a 40-cell factorial at the prazosin-PTSD reference design
with 1000 replicates per power cell and 2000 per Type I cell.
The continuous-biomarker contrast differs from strict RM-ANOVA
only in retaining the continuous biomarker, isolating the cost
of dichotomization. The cycle-by-period design grid is specified
here but its empirical sweep is deferred to a companion paper.

\textbf{Results.} The closed-form strict RM-ANOVA test agrees with the
linear-mixed Wald test under sphericity and balanced design. In
the simulation, most of the power gap between strict RM-ANOVA and
the mixed model is the cost of dichotomization rather than of the
test procedure: at \(N = 25\), replacing the dichotomized biomarker
with the continuous one raises power by 0.180, while the further
step to the full mixed model adds only 0.027, a pattern stable
across \(N\) (dichotomization 0.08--0.21, test procedure at most
0.04). Among the four continuous-biomarker procedures, the
linear-mixed analysis attains the highest power within nominal
Type I error; GEE with the naive sandwich inflates Type I error
to between \(1.5\times\) and \(1.9\times\) nominal across all four
sample sizes, which the Mancl-DeRouen correction removes at a
moderate (about 10-percentage-point) power cost while restoring
93--95\% interval coverage. The inline Mancl-DeRouen standard
error agrees with the \texttt{geesmv} reference to within 2.5\%.

\textbf{Conclusions.} The analyst's most consequential decision is to
keep the biomarker continuous rather than dichotomize it; among
continuous-biomarker procedures the linear-mixed analysis is the
preferred default and a dichotomization-free within-subject
contrast is a close, transparent alternative. Naive GEE should
not be used at the sample sizes examined without a small-sample
correction. This paper resolves the test-procedure axis at a
fixed reference design; the joint cycle-by-period design grid is
specified here and deferred to a companion paper.

\section{Introduction}\label{introduction}

\subsection{A claim about how design and test choices interact}\label{a-claim-about-how-design-and-test-choices-interact}

The present paper develops the methodological case for treating
the analyst's two principal choices in an aggregated N-of-1 trial
jointly rather than separately. The two choices are the test
procedure for the biomarker-treatment interaction (strict
repeated-measures ANOVA, the linear mixed-effects model with
continuous-time AR(1) residual correlation, the mixed-effects
model for repeated measures, generalized estimating equations)
and the trial-design parameters that govern within-subject
treatment cycling (number of treatment cycles, on-drug and
off-drug period lengths, on/off symmetry, optional run-in and
run-out phases). The validity and power of any given test depend
on what the design assumes about carryover, autocorrelation, and
within-subject variance structure; conversely, the optimal design
parameters depend on which test the analyst plans to use.
Treating the two as independent optimizations, as the published
methodological literature typically does, produces locally optimal
but jointly suboptimal combinations, and we shall argue that this
gap is consequential at the sample sizes characteristic of
precision-medicine pilot studies.

\subsection{Why the question matters for chronic-condition pharmacology}\label{why-the-question-matters-for-chronic-condition-pharmacology}

In this section we shall set out three reasons why the joint
treatment of test procedure and design parameters is warranted
specifically for the aggregated N-of-1 design class.

We begin with the physiological motivation. Within-participant
residuals in chronic-condition trials are autocorrelated not as
an analytic convenience but because most chronic-condition
response trajectories exhibit serial correlation on a timescale
set by drug pharmacokinetics, condition-specific natural-history
drift, and measurement reactivity (Hawthorne, structured-visit,
and diary-completion effects). Continuous-time AR(1) residual
correlation is, in our view, the appropriate structural prior for
this physiology, but it does not satisfy the sphericity assumption
required by strict RM-ANOVA.\textsuperscript{\citeproc{ref-haverkamp2017sphericity}{2}} show
through simulation that uncorrected RM-ANOVA inflates Type I error
under sphericity violations; they also report, importantly, that
the Greenhouse-Geisser-corrected RM-ANOVA and the multilevel
linear model perform comparably once the correction is applied,
and that the multilevel model is not uniformly safer (it can be
the more liberal of the two in some conditions). The
degrees-of-freedom corrections\textsuperscript{\citeproc{ref-huynh1976}{3},\citeproc{ref-greenhouse1959}{4}}
recover nominal levels only conditionally.\textsuperscript{\citeproc{ref-blanca2023sphericity}{5}}
extend this evidence with finer epsilon resolution, with the
conventional rule favoring Greenhouse-Geisser when
\(\epsilon < 0.75\). The advantage we claim for the mixed-effects
model with explicit \texttt{corCAR1} residual structure is thus not that
it is uniformly more powerful than a well-corrected RM-ANOVA, but
that it models the AR(1) structure directly, accommodates the
continuous biomarker and unequal time spacing, and so avoids both
the dichotomization and the degrees-of-freedom penalty that the
classical route incurs.

\subsection{What the published N-of-1 methodology literature does and does not do}\label{what-the-published-n-of-1-methodology-literature-does-and-does-not-do}

We begin with the state of the analytic-model literature. The
systematic review of\textsuperscript{\citeproc{ref-natesanbatley2023}{6}} documents that 83.8\% of
the published N-of-1 studies they catalogue ignore autocorrelation
entirely;\textsuperscript{\citeproc{ref-stunnenberg2022}{7}} and\textsuperscript{\citeproc{ref-shaffer2018}{8}} reach similar
conclusions in neurology and behavioral-medicine subdomains.
CONSORT extension reporting standards\textsuperscript{\citeproc{ref-vohra2015cent}{9},\citeproc{ref-shamseer2015cent}{10}} codify reporting quality but do not address
the analytic-model decision in any detail, and\textsuperscript{\citeproc{ref-kwasnicka2019}{11}}
identifies the analytic-method gap as the highest-priority
methodological problem in the field. It appears, then, that the
majority of N-of-1 analyses are conducted under a misspecified
residual covariance assumption, and that the field has not yet
converged on remedies.

Test-procedure comparisons for longitudinal data are not new. The
generic contrast of RM-ANOVA against GEE and mixed-effects models
is well established:\textsuperscript{\citeproc{ref-ma2012beyond}{12}}, for instance, report through
simulation that RM-ANOVA is appreciably less powerful than GEE
and mixed models for continuous longitudinal endpoints, the same
ordering we recover. What that literature does not do is treat
the biomarker-treatment interaction estimand, the aggregated
N-of-1 / cross-over design, or the continuous-time AR(1)
structure together, nor does it isolate the dichotomization cost
that the classical route imposes specifically when the moderator
is continuous. The MMRM consensus established by Mallinckrodt and
colleagues\textsuperscript{\citeproc{ref-mallinckrodt2003}{13}--\citeproc{ref-prakash2008}{16}} applies to
parallel-group longitudinal trials with monotone missingness, not
to aggregated N-of-1 designs. The GEE methodology of\textsuperscript{\citeproc{ref-liang1986}{17}}
with sandwich variance estimation has been adapted to
small-cluster settings via bias-corrected estimators\textsuperscript{\citeproc{ref-wang2016geesmv}{18},\citeproc{ref-pan2001}{19}}, but the small-sample concerns are
acute at trial-relevant N-of-1 cluster counts:\textsuperscript{\citeproc{ref-tong2024stepwedge}{20}}
report that in stepped-wedge cluster trials with fewer than six
independent units only 6.6\% of analyses use appropriate
small-sample corrections, illustrating the gap between method
availability and practice. None of these papers address the
biomarker-interaction estimand in the cross-over or N-of-1
setting directly.

The gap, then, is this. Across the recent systematic reviews\textsuperscript{\citeproc{ref-natesanbatley2023}{6},\citeproc{ref-stunnenberg2022}{7},\citeproc{ref-defelippe2023}{21}} and the
methodological literature surveyed above, no published study
crosses test procedure (strict RM-ANOVA, MMRM, GEE,
mixed-effects with \texttt{corCAR1}) against a design-parameter grid
(cycles, period lengths, asymmetry, run-in, run-out) under
simulation for the aggregated N-of-1 design class. The present
paper endeavors to address that gap, at least on the
test-procedure axis.

\subsection{Run-in design and the placebo-selection controversy}\label{run-in-design-and-the-placebo-selection-controversy}

\subsection{What this paper contributes}\label{what-this-paper-contributes}

We make four contributions, each addressed in a section below.

First, we derive the closed-form strict RM-ANOVA \(F\)-statistic
for the biomarker-stratified treatment-by-time interaction in a
balanced split-plot design with categorical predictors and the
sphericity assumption (\S2). We relate the strict-RM-ANOVA test
to the Wald \(t\)-test on the biomarker-treatment
fixed-effects coefficient in a linear mixed model with
continuous-time AR(1) residual correlation, identifying the
conditions on residual covariance under which the two tests agree
and the regimes in which they diverge. The derivation extends
standard cross-over methodology\textsuperscript{\citeproc{ref-senn1994crossover}{22},\citeproc{ref-senn2002crossover}{23}} to the biomarker-stratified case and clarifies
the connection to the contemporary sphericity-violation literature\textsuperscript{\citeproc{ref-haverkamp2017sphericity}{2},\citeproc{ref-blanca2023sphericity}{5}}.

Third, we conduct a Monte Carlo simulation study (\S4--\S5) that
crosses five test procedures (strict RM-ANOVA on the dichotomized
biomarker; a continuous-biomarker within-subject contrast; the
standard pmsimstats linear-mixed analysis with \texttt{corCAR1}; and GEE
with naive and with Mancl-DeRouen\textsuperscript{\citeproc{ref-mancl2001covariance}{24}}
small-sample sandwich variance) against the biomarker-moderation
coefficient and sample size at the fixed reference design. The
continuous-biomarker contrast is the dichotomization-free
analogue of the strict RM-ANOVA test and lets us separate the
cost of dichotomizing the biomarker from the cost of the test
procedure proper. The simulation runs on the pmsimstats package
data-generating process under Architecture A (mean moderation)
via \texttt{generateData()} and \texttt{lme\_analysis()}, holding fixed the
response-curve and covariance-structure axes addressed in
\texttt{analysis/report/07-gompertz-evaluation/} and
\texttt{analysis/report/01-dgp-mean-moderation-vs-mvn/}. Results are
reported in the\textsuperscript{\citeproc{ref-morris2019}{1}} ADEMP framework adopted as the
project-wide simulation-reporting standard.

\section{Closed-form strict RM-ANOVA test of the biomarker-treatment interaction}\label{closed-form-strict-rm-anova-test-of-the-biomarker-treatment-interaction}

In this section we shall derive the closed-form strict RM-ANOVA
\(F\)-statistic for the biomarker-treatment interaction in a
balanced split-plot design with categorical predictors and the
sphericity assumption, and relate it to the Wald \(t\)-statistic
on the corresponding linear-mixed coefficient. The derivation
specializes classical split-plot machinery\textsuperscript{\citeproc{ref-maxwelldelaney2018}{25}}
to the biomarker-stratified case in the simplest design that
supports the test: two biomarker strata, two treatment phases,
and \(T\) within-phase timepoints. The presentation is deliberately
confined to strict RM-ANOVA. Modern extensions --- the
mixed-effects model for repeated measures\textsuperscript{\citeproc{ref-mallinckrodt2008recommendations}{26}}, generalized estimating
equations\textsuperscript{\citeproc{ref-liang1986}{17}}, and ANCOVA with random subject effects ---
appear only for context; the closed-form expression we report is
the textbook compound-symmetric form. Three findings anchor the
rest of the paper:

\begin{enumerate}
\def\labelenumi{\arabic{enumi}.}
\tightlist
\item
  Under sphericity and balance, the strict RM-ANOVA test of
  the biomarker-treatment interaction is identical to a
  two-sample \(t\)-test on the within-subject treatment
  differences, with degrees of freedom \(2(n-1)\).
\item
  The closed form yields a clean non-central \(F\) power formula
  that we use as the analytic anchor for the simulation study
  in §4-§5.
\item
  Five assumptions of the closed form are violated by the
  pmsimstats data-generating process (DGP; continuous biomarker,
  AR(1) within-subject covariance, exponential carryover,
  missing data, unequal time spacing). Each violation
  favors the linear-mixed analysis over the strict RM-ANOVA
  test in operating characteristics, motivating the
  primary-and-sensitivity reporting strategy of §6.
\end{enumerate}

\subsection{Notation and design}\label{notation-and-design}

Specifically, we consider the following structural elements:

\begin{itemize}
\tightlist
\item
  Two biomarker strata \(i \in \{1, 2\}\), formed by dichotomizing
  the continuous biomarker at a pre-specified cut-point. Strict
  RM-ANOVA requires categorical predictors; the
  dichotomization is the price of admission to the framework.
\item
  \(n\) subjects per stratum, balanced. Index
  \(l \in \{1, \ldots, n\}\).
\item
  Two treatment phases \(j \in \{1, 2\}\), active drug and placebo.
  In an N-of-1 design these are two phases through which every
  subject passes; in a parallel-group design they are
  between-subject and the design is not a split-plot.
\item
  \(T\) within-phase timepoints, indexed
  \(k \in \{1, \ldots, T\}\). In the simplest case \(T = 1\), the
  analyst pre-averages within-phase observations to a single
  mean per subject per phase.
\end{itemize}

We write \(Y_{ijlk}\) for the symptom outcome at biomarker stratum
\(i\), treatment phase \(j\), subject \(l\) within stratum \(i\), and
timepoint \(k\). The strict classical split-plot model is

\begin{equation}
\begin{split}
Y_{ijlk} = {} & \mu + \alpha_i + \beta_j + (\alpha\beta)_{ij}
+ \pi_{l(i)} \\
& + \tau_k + (\alpha\tau)_{ik} + (\beta\tau)_{jk}
+ (\alpha\beta\tau)_{ijk} + e_{ijlk},
\end{split}
\label{eq:rmanova-model}
\end{equation}

with sum-to-zero constraints on each fixed effect, the random
subject effect \(\pi_{l(i)} \sim N(0, \sigma_\pi^2)\) nested in
stratum, and within-subject error \(e_{ijlk} \sim N(0,
\sigma_e^2)\). The model presupposes the classical sphericity
(compound symmetry) condition on the within-subject covariance:
every pair of within-subject observations has the same covariance
\(\sigma_\pi^2\), and every observation has the same total variance
\(\sigma_\pi^2 + \sigma_e^2\).

\subsection{\texorpdfstring{Sums of squares and the interaction \(F\)-statistic}{Sums of squares and the interaction F-statistic}}\label{sums-of-squares-and-the-interaction-f-statistic}

We begin the derivation by partitioning the total sum of squares.
The total sum of squares partitions into between-subjects and
within-subjects components:

\begin{equation}
SS_{\text{tot}} = SS_{\text{between}} + SS_{\text{within}} .
\label{eq:rmanova-ss-partition}
\end{equation}

The between-subjects component decomposes further into \(SS_A\)
(biomarker stratum) and \(SS_{S(A)}\) (subject within stratum):

\begin{equation}
SS_A = bnT \sum_{i=1}^a (\bar Y_{i\cdots} - \bar Y_{\cdots})^2,
\quad
SS_{S(A)} = bT \sum_{i=1}^a \sum_{l=1}^n
(\bar Y_{i\cdot l\cdot} - \bar Y_{i\cdots})^2 .
\label{eq:rmanova-between}
\end{equation}

The within-subjects component decomposes into \(SS_B\) (treatment
phase), \(SS_{AB}\) (the biomarker-treatment interaction, the
primary target of the test), and a within-subjects error term
\(SS_{B \times S(A)}\):

\begin{equation}
SS_{AB} = nT \sum_{i,j} (\bar Y_{ij\cdot\cdot}
- \bar Y_{i\cdots} - \bar Y_{\cdot j\cdot\cdot}
+ \bar Y_{\cdots})^2 ,
\label{eq:rmanova-ss-ab}
\end{equation}

\begin{equation}
SS_{B \times S(A)} = T \sum_{i,j,l} (\bar Y_{ijl\cdot}
- \bar Y_{i\cdot l\cdot} - \bar Y_{ij\cdot\cdot}
+ \bar Y_{i\cdots})^2 .
\label{eq:rmanova-ss-error}
\end{equation}

Corresponding degrees of freedom are \(df_A = a - 1\),
\(df_{S(A)} = a(n-1)\), \(df_B = b - 1\), \(df_{AB} = (a-1)(b-1)\),
and \(df_{B \times S(A)} = a(n-1)(b-1)\), the last being the
within-subjects error degrees of freedom appropriate for the
treatment effect and its interaction with the between-subjects
factor. The biomarker-treatment interaction is then tested by the
\(F\)-ratio

\begin{equation}
F_{AB} = \frac{MS_{AB}}{MS_{B \times S(A)}}
= \frac{SS_{AB} / df_{AB}}
       {SS_{B \times S(A)} / df_{B \times S(A)}} .
\label{eq:rmanova-f}
\end{equation}

Under the null \(H_0: (\alpha\beta)_{ij} \equiv 0\) and the
sphericity assumption, \(F_{AB} \sim F(df_{AB}, df_{B \times
S(A)})\).

\subsubsection{\texorpdfstring{The \(2 \times 2\) split-plot case}{The 2 \textbackslash times 2 split-plot case}}\label{the-2-times-2-split-plot-case}

We now specialize to \(a = b = 2\) (two biomarker strata, two
phases), setting \(T = 1\) for cleanest exposition; that is, the
analyst pre-averages within-phase observations to a single value
per subject per phase before computing the interaction. The
sample interaction contrast is

\begin{equation}
\hat L = (\bar Y_{11\cdot} - \bar Y_{12\cdot})
- (\bar Y_{21\cdot} - \bar Y_{22\cdot}) ,
\label{eq:rmanova-Lhat}
\end{equation}

That is, \(\hat L\) is the difference of within-subject treatment
differences across biomarker strata, and is the natural
biomarker-treatment interaction statistic. The interaction
parameter under the sum-to-zero ANOVA parameterization satisfies
\(\widehat{(\alpha\beta)}_{11} = \hat L / 4\) (and the other three
cells take values \(\pm \hat L / 4\)), so the interaction sum of
squares is \(SS_{AB} = nT\,\hat L^2 / 4\), and with \(df_{AB} = 1\)
this equals the interaction mean square. Under sphericity and
balance, \(MS_{B \times S(A)} = \hat\sigma_e^2\). Substituting
into equation \eqref{eq:rmanova-f} yields

\begin{equation}
F_{AB} = \frac{nT\,\hat L^2}{4\,\hat\sigma_e^2},
\qquad df_1 = 1, \quad df_2 = 2(n-1) .
\label{eq:rmanova-f-22}
\end{equation}

\subsubsection{\texorpdfstring{Equivalence to a within-subject \(t\)-test}{Equivalence to a within-subject t-test}}\label{equivalence-to-a-within-subject-t-test}

We note that equation \eqref{eq:rmanova-f-22} is equivalent to a
two-sample \(t\)-test on the within-subject differences. Define
\(\Delta_{il} = \bar Y_{i1l\cdot} - \bar Y_{i2l\cdot}\) as the
within-subject treatment difference for subject \(l\) in biomarker
stratum \(i\), averaged over \(T\) within-phase timepoints. Then

\begin{equation}
\hat L = \bar\Delta_1 - \bar\Delta_2,
\qquad
\bar\Delta_i = \frac{1}{n}\sum_{l=1}^n \Delta_{il}.
\label{eq:rmanova-Delta}
\end{equation}

Under sphericity, \(\Delta_{il}\) are i.i.d.~within stratum with
\(\mathrm{Var}(\Delta_{il}) = 2\sigma_e^2 / T\), since the random
subject effect cancels in the within-subject difference. The
variance of \(\hat L\) is therefore

\begin{equation}
\mathrm{Var}(\hat L) = \frac{2 \cdot 2\sigma_e^2}{nT}
= \frac{4\sigma_e^2}{nT} ,
\label{eq:rmanova-varL}
\end{equation}

and the \(t\)-statistic for the contrast is

\begin{equation}
t_{AB} = \frac{\hat L}{\sqrt{4\hat\sigma_e^2 / (nT)}}
= \sqrt{\frac{nT}{4}} \cdot \frac{\hat L}{\hat\sigma_e},
\label{eq:rmanova-t}
\end{equation}

\subsection{Relation to the linear-mixed Wald test}\label{relation-to-the-linear-mixed-wald-test}

We shall now relate the strict-RM-ANOVA \(F\)-statistic of
equation \eqref{eq:rmanova-f-22} to the pmsimstats \texttt{nlme::lme}
Wald \(t\)-statistic on the \texttt{bm:Dbc} coefficient. The two are
asymptotically equivalent under three conditions: (i) the
biomarker is dichotomized at the cut-point that defines the
strata, (ii) the within-subject covariance is exactly
compound-symmetric, and (iii) \texttt{Dbc} is binary, the A1
specification of the carryover-sensitivity manuscript rather than
the matched-decay A2 specification. When all three conditions
hold, the two tests reach the same coefficient and the same
standard error up to small-sample correction factors. When any
one fails, the mixed-model test extracts more information than
the strict RM-ANOVA: the continuous biomarker contributes
within-stratum variation, the AR(1) covariance is modeled
directly, and the continuous exposure-decay predictor draws on
the off-drug trajectory.

The pmsimstats DGP imposes AR(1) within-factor correlation, which
is not compound-symmetric. Under AR(1), Box's \(\epsilon\) is
strictly below 1, and treating the full multi-timepoint series as
repeated within-subject levels under classical RM-ANOVA inflates
Type I error unless a degrees-of-freedom correction is applied.
The relevant magnitude is recorded in the project's own audit
(\texttt{docs/06-ar1-residual-correlation.tex}): at \(T = 8\) timepoints,
fitting a mixed model with no residual-correlation structure (the
naive analogue of treating AR(1) data as exchangeable) produces
Type I error of 13 to 17 percent for the crossover (CO) and
open-label (OL) designs, which switching to \texttt{nlme::lme} with
\texttt{corCAR1} returns to nominal. We stress that this 13-to-17-percent
figure is for the unstructured/naive fit, not for the strict
RM-ANOVA arm evaluated in this paper: that arm pre-averages each
phase to a single value (the \(2 \times 2\) form of \S2.2), so its
within-subject covariance is \(2 \times 2\), sphericity holds by
construction, and its Type I error is at nominal (\S5). The
sphericity penalty is therefore a property of the un-pre-averaged
classical analysis; for the pre-averaged strict RM-ANOVA the
operative cost is dichotomization, not sphericity. Either way the
linear-mixed analysis with \texttt{corCAR1} avoids both costs, modeling
the AR(1) structure directly rather than absorbing it via a
degrees-of-freedom penalty.

\subsubsection{Power calculation under the closed form}\label{power-calculation-under-the-closed-form}

For the \(2 \times 2\) case, equation \eqref{eq:rmanova-f-22} yields
a closed-form power formula that we shall use as the analytic
anchor for the simulation study in \S4--\S5. Under the
alternative \(H_1: \hat L = \delta\), the statistic
\(F_{AB} = nT\,\hat L^2 / (4\hat\sigma_e^2)\) follows a
non-central \(F\) distribution with non-centrality parameter

\begin{equation}
\lambda = \frac{nT\,\delta^2}{4\sigma_e^2} ,
\label{eq:rmanova-ncp}
\end{equation}

and power at level \(\alpha\) is

\begin{equation}
1 - \beta = \Pr\!\left[\, F_{AB} > F_{1, 2(n-1), 1-\alpha}
\,\big|\, \lambda \,\right] ,
\label{eq:rmanova-power}
\end{equation}

\subsubsection{Limitations of the strict classical formulation}\label{limitations-of-the-strict-classical-formulation}

Unfortunately, the closed-form simplicity of equation
\eqref{eq:rmanova-f-22} is bought at the price of five strong
assumptions, each of which is empirically problematic in the
pmsimstats setting. We present them in approximate order of
severity.

\begin{enumerate}
\def\labelenumi{\arabic{enumi}.}
\item
  \emph{Categorical biomarker.} The pmsimstats biomarker (resting
  SBP in the prazosin-PTSD application) is continuous and has
  no defensible cut-point. Dichotomizing at the median or at
  a clinical threshold loses the within-stratum biomarker
  variation, attenuates the interaction estimator, and
  inflates Type II error. In the carryover-sensitivity
  production data the median split gives a \(\hat L\) standard
  error roughly 1.4 to 1.7 times that of the continuous-
  biomarker linear-mixed coefficient, translating to a
  30 to 50 percent power penalty.
\item
  \emph{Sphericity.} Compound-symmetric within-subject covariance
  is empirically rejected for time-series symptom data:
  nearby timepoints are more correlated than distant ones.
  The AR(1) structure pmsimstats implements explicitly
  violates sphericity. After Greenhouse-Geisser correction,
  the test recovers nominal Type I error but loses degrees
  of freedom and power.
\item
  \emph{No carryover model.} Strict RM-ANOVA has no provision for
  residual drug effect during the off-drug period; the model treats
  off-drug observations as if no drug had been administered.
  Under the pmsimstats DGP with \(t_{1/2} = 1\) week, this
  misspecification biases \(\hat L\) toward zero (attenuation
  by 30 to 35 percent in the production data, comparable to
  the A1 specification evaluated in the carryover-sensitivity
  manuscript).
\item
  \emph{Balanced design required.} Equations
  \eqref{eq:rmanova-f} and \eqref{eq:rmanova-f-22} assume equal
  stratum sizes and equal numbers of within-phase timepoints.
  Real trials have missing data; classical RM-ANOVA handles
  missingness only through complete-case deletion. The MMRM
  framework's likelihood-based handling of missing-at-random (MAR)
  data\textsuperscript{\citeproc{ref-mallinckrodt2008recommendations}{26}} is genuinely
  better-behaved.
\item
  \emph{Equal time spacing.} Strict RM-ANOVA treats time as a
  categorical factor with implicitly equal-spaced levels.
  Real N-of-1 trials have unequally-spaced visits; pmsimstats
  implements continuous-time AR(1) (\texttt{corCAR1}), which the
  strict-RM-ANOVA framework cannot accommodate.
\end{enumerate}

These limitations motivate retaining the linear-mixed analysis
with continuous-time AR(1) residual correlation as primary, and
reporting the strict-RM-ANOVA closed-form \(F\)-statistic as a
sensitivity comparator. The five violations are not equally
consequential: the categorical-biomarker and no-carryover
assumptions are the most problematic in the pmsimstats setting,
together accounting for a 30 to 50 percent attenuation of the
interaction estimator relative to the linear-mixed analysis. The
recommendation operationalized in \S6 follows this strategy.

\section{Cycle-by-period design grid}\label{cycle-by-period-design-grid}

The cycle-by-period design grid that varies the number of
treatment cycles, on-drug and off-drug period lengths, and on/off
symmetry around the\textsuperscript{\citeproc{ref-hendrickson2020}{27}} hybrid reference is
specified in \texttt{docs/27-cycle-period-design-sweep-plan.md} and
constitutes the second axis of the joint design-by-test
sensitivity question framed in \S1. The empirical sweep across
this grid is deferred to a follow-up paper. We shall note here,
in the interest of transparency, that the present paper narrows
its scope to the test-procedure axis at the prazosin/PTSD
reference design (Hendrickson hybrid open-label plus blinded
discontinuation, OL+BDC; \(N \in \{25, 35, 70, 100\}\),
\(c_{bm} \in \{0, 0.45\}\)), which the simulation infrastructure
does support at production rep counts. The \S6 recommendation is
correspondingly scoped to test-procedure choice rather than the
joint (cycle count, period length, test procedure) optimum that
the design-by-test sensitivity question requires. The follow-up paper executing the
cycle-by-period sweep is the natural extension and is recorded as
a portfolio item.

\section{Simulation design}\label{simulation-design}

The simulation program follows the ADEMP framework of\textsuperscript{\citeproc{ref-morris2019}{1}}. The primary estimand is the biomarker-treatment
interaction inference under each of the five test procedures.
Four operate on the continuous biomarker and share the same
interaction estimand \(\beta_{bm:D}\) on the symptom scale: a
within-subject treatment-effect contrast regressed on the
continuous biomarker (the dichotomization-free analogue of the
two-sample \(t\)-test of equation \eqref{eq:rmanova-t}); the
linear-mixed-effects Wald \(t\) on \(\hat{\beta}_{bm:D}\); the GEE
Wald \(z\) on \(\hat{\beta}_{bm:D}\) with naive sandwich variance;
and the GEE Wald \(z\) with Mancl-DeRouen small-sample sandwich
variance per\textsuperscript{\citeproc{ref-mancl2001covariance}{24}}. The fifth, the strict
classical RM-ANOVA \(F\)-statistic, operates on the
median-dichotomized biomarker and so targets a dichotomized
estimand on a different scale. The continuous-biomarker contrast
is included specifically to separate the cost of dichotomization
(strict RM-ANOVA versus the contrast, both pre-averaged to one
value per phase) from the cost of the test procedure proper (the
contrast versus the longitudinal linear-mixed and GEE analyses).
Primary inferential outputs are power \(\pi = \Pr(p < 0.05)\) for
\(H_0: \beta_{bm:D} = 0\) at \(\alpha = 0.05\), and Type I error
under \(\beta_{bm:D} = 0\). Secondary estimands, reported for the
four coefficient-bearing arms, are bias of \(\hat{\beta}_{bm:D}\)
relative to a high-precision Monte-Carlo reference value of the
estimand, the empirical and the mean model-based standard error,
and nominal-95\% CI coverage. Strict RM-ANOVA is excluded from
these coefficient-based measures because its estimand is on the
dichotomized scale.

The full Aims, Data-generating mechanisms, Estimands, Methods,
and Performance measures are pre-registered at
\texttt{analysis/scripts/test-procedure-design-sensitivity/00-ademp-pre-registration.md},
which fixes the test-procedure times design-grid factorial across
three studies (Study 1: test-procedure comparison at fixed
design; Study 2: cycle-by-period grid at fixed test, 240 cells;
Study 3: joint design-by-test heatmap, 240 cells) before the
production runs commit. The present paper executes a revised
Study 1; the executed design departs from the pre-registered
Study 1 on several factor levels, and we document every
departure in Table \ref{tab:deviations} below in the interest
of full ADEMP transparency.

\begin{longtable}[]{@{}
  >{\raggedright\arraybackslash}p{(\linewidth - 4\tabcolsep) * \real{0.2100}}
  >{\raggedright\arraybackslash}p{(\linewidth - 4\tabcolsep) * \real{0.3100}}
  >{\raggedright\arraybackslash}p{(\linewidth - 4\tabcolsep) * \real{0.4800}}@{}}
\caption{\label{tab:deviations} Pre-registered versus executed Study 1.
Departures are recorded for transparency; the executed superset
of sample sizes (which restores the pre-registered \(N\) levels
and adds \(N = 25\)) and the added continuous-biomarker contrast
arm respond to referee comments on the first internal draft.}\tabularnewline
\toprule\noalign{}
\begin{minipage}[b]{\linewidth}\raggedright
Factor
\end{minipage} & \begin{minipage}[b]{\linewidth}\raggedright
Pre-registered
\end{minipage} & \begin{minipage}[b]{\linewidth}\raggedright
Executed (this paper)
\end{minipage} \\
\midrule\noalign{}
\endfirsthead
\toprule\noalign{}
\begin{minipage}[b]{\linewidth}\raggedright
Factor
\end{minipage} & \begin{minipage}[b]{\linewidth}\raggedright
Pre-registered
\end{minipage} & \begin{minipage}[b]{\linewidth}\raggedright
Executed (this paper)
\end{minipage} \\
\midrule\noalign{}
\endhead
\bottomrule\noalign{}
\endlastfoot
Test procedures & 3 (RM-ANOVA, LME, GEE-MD) & 5 (adds naive GEE and the continuous-biomarker contrast) \\
Biomarker effect \(c_{bm}\) & \(\{0, 0.3, 0.6\}\) & \(\{0, 0.45\}\) (the project-wide reference effect size) \\
Sample size \(N\) & \(\{35, 70, 100\}\) & \(\{25, 35, 70, 100\}\) (superset; adds \(N = 25\)) \\
Replicates per cell & 1000 (5000 for Type I cells) & 1000 power, 2000 Type I \\
Carryover half-life & \(t_{1/2} = 3\) days & \(t_{1/2} = 0\) (no carryover at the fixed reference) \\
\end{longtable}

In the scope of the present paper (test-procedure axis only at
the reference design), the simulation runs on the pmsimstats
package data-generating process under Architecture A (mean
moderation), exercised through \texttt{generateData()} and the standard
\texttt{lme\_analysis()} fit, rather than on the reduced
\texttt{implementations/simple} substrate. The factorial comprises
\(n_{\text{reps}} = 1000\) per power cell and \(2000\) per Type I
cell, across the 40 cells defined by procedure \(\in \{\)strict
RM-ANOVA, continuous-biomarker contrast, linear-mixed with
\texttt{corCAR1}, GEE with naive sandwich, GEE with Mancl-DeRouen
small-sample correction\(\} \times c_{bm} \in \{0, 0.45\} \times
N \in \{25, 35, 70, 100\}\). Execution uses
\texttt{parallel::mclapply} with \texttt{RNGkind("L\textquotesingle{}Ecuyer-CMRG")} and
\texttt{mc.set.seed\ =\ TRUE}, so each replicate draws from an
independent, reproducible random-number stream. The
Mancl-DeRouen correction was implemented inline (no
\texttt{geesmv}-style external dependency in the production path) by
inflating each cluster's residual via \((I - H_{ii})^{-1}\,r_i\)
in the sandwich meat. It produces identical point estimates and
larger standard errors than the naive sandwich, and its SE for
the interaction coefficient agrees with \texttt{geesmv::GEE.var.md} to
within approximately 2.5\% (mean inline/\texttt{geesmv} ratio 1.02,
standard deviation 0.02, over 25 probe replicates at
\(N = 35\), \(c_{bm} = 0.45\)), so the correction magnitude is
externally validated rather than merely plausible.

\section{Results}\label{results}

The simulation program was executed across the 40 cells of the
five-procedure-by-effect-by-\(N\) factorial in 587 s using 7-core
\texttt{parallel::mclapply} with \texttt{RNGkind("L\textquotesingle{}Ecuyer-CMRG")} per-stream
seeding. Convergence was 100\% across all 60\{,\}000 fits.
Per-replicate output is at
\texttt{analysis/data/quick-sim/08-test-replicates.rds}; the Morris
ADEMP cell-level summary is at
\texttt{analysis/data/quick-sim/08-test-summary.txt}. With 1000
replicates per power cell the MCSE on a power estimate near
\(0.5\) is approximately \(0.016\), and with 2000 replicates per
Type I cell the MCSE on a rate near \(0.05\) is approximately
\(0.005\). We begin the results with the test-procedure comparison
at the fixed reference design.

\subsection{Test-procedure comparison at fixed design}\label{test-procedure-comparison-at-fixed-design}

We note at the outset that strict RM-ANOVA targets a different
estimand from the other four procedures. It operates on a
median-dichotomized biomarker, pre-averaged to one value per
phase, and so estimates a dichotomized-stratum interaction; the
continuous-biomarker contrast, the linear-mixed model, and both
GEE variants estimate the interaction coefficient
\(\beta_{bm:D}\) on the continuous symptom scale. An earlier
version of this comparison contrasted strict RM-ANOVA directly
with the mixed model and attributed the resulting power gap to
the test procedure; that attribution conflated two distinct
costs. We therefore introduce the continuous-biomarker contrast
of \S4, which differs from strict RM-ANOVA only in retaining the
continuous biomarker (it shares the pre-averaging and the absence
of any AR(1) machinery) and differs from the linear-mixed model
only in discarding the within-phase longitudinal structure. The
RM-ANOVA \(\to\) contrast step therefore measures the cost of
dichotomization in isolation, and the contrast \(\to\) linear-mixed
step measures the cost of the test procedure proper. Reading the
two steps separately is the honest decomposition the earlier
direct comparison lacked.

Type I error at the null cell (\(c_{bm} = 0\)), percent \(\pm\)
MCSE, is shown in Table \ref{tab:type1}; power at
\(c_{bm} = 0.45\) in Table \ref{tab:power}.

\begin{longtable}[]{@{}
  >{\raggedright\arraybackslash}p{(\linewidth - 8\tabcolsep) * \real{0.2917}}
  >{\raggedright\arraybackslash}p{(\linewidth - 8\tabcolsep) * \real{0.1771}}
  >{\raggedright\arraybackslash}p{(\linewidth - 8\tabcolsep) * \real{0.1771}}
  >{\raggedright\arraybackslash}p{(\linewidth - 8\tabcolsep) * \real{0.1771}}
  >{\raggedright\arraybackslash}p{(\linewidth - 8\tabcolsep) * \real{0.1771}}@{}}
\caption{\label{tab:type1} Type I error rate (percent \(\pm\) MCSE) under
\(H_0: \beta_{bm:D} = 0\), 2000 replicates per cell. Nominal level
5\%.}\tabularnewline
\toprule\noalign{}
\begin{minipage}[b]{\linewidth}\raggedright
Procedure
\end{minipage} & \begin{minipage}[b]{\linewidth}\raggedright
\(N = 25\)
\end{minipage} & \begin{minipage}[b]{\linewidth}\raggedright
\(N = 35\)
\end{minipage} & \begin{minipage}[b]{\linewidth}\raggedright
\(N = 70\)
\end{minipage} & \begin{minipage}[b]{\linewidth}\raggedright
\(N = 100\)
\end{minipage} \\
\midrule\noalign{}
\endfirsthead
\toprule\noalign{}
\begin{minipage}[b]{\linewidth}\raggedright
Procedure
\end{minipage} & \begin{minipage}[b]{\linewidth}\raggedright
\(N = 25\)
\end{minipage} & \begin{minipage}[b]{\linewidth}\raggedright
\(N = 35\)
\end{minipage} & \begin{minipage}[b]{\linewidth}\raggedright
\(N = 70\)
\end{minipage} & \begin{minipage}[b]{\linewidth}\raggedright
\(N = 100\)
\end{minipage} \\
\midrule\noalign{}
\endhead
\bottomrule\noalign{}
\endlastfoot
Strict RM-ANOVA (dich.) & 4.55 \(\pm\) 0.47 & 4.15 \(\pm\) 0.45 & 5.95 \(\pm\) 0.53 & 5.20 \(\pm\) 0.50 \\
Contrast (continuous bm) & 5.25 \(\pm\) 0.50 & 4.50 \(\pm\) 0.46 & 6.45 \(\pm\) 0.55 & 4.60 \(\pm\) 0.47 \\
Linear-mixed (\texttt{corCAR1}) & 2.75 \(\pm\) 0.37 & 2.05 \(\pm\) 0.32 & 3.35 \(\pm\) 0.40 & 2.90 \(\pm\) 0.38 \\
GEE (naive sandwich) & 9.70 \(\pm\) 0.66 & 8.55 \(\pm\) 0.63 & 7.90 \(\pm\) 0.60 & 7.40 \(\pm\) 0.59 \\
GEE (Mancl-DeRouen) & 6.35 \(\pm\) 0.55 & 5.30 \(\pm\) 0.50 & 6.25 \(\pm\) 0.54 & 5.90 \(\pm\) 0.53 \\
\end{longtable}

\begin{longtable}[]{@{}
  >{\raggedright\arraybackslash}p{(\linewidth - 8\tabcolsep) * \real{0.2800}}
  >{\raggedright\arraybackslash}p{(\linewidth - 8\tabcolsep) * \real{0.1800}}
  >{\raggedright\arraybackslash}p{(\linewidth - 8\tabcolsep) * \real{0.1800}}
  >{\raggedright\arraybackslash}p{(\linewidth - 8\tabcolsep) * \real{0.1800}}
  >{\raggedright\arraybackslash}p{(\linewidth - 8\tabcolsep) * \real{0.1800}}@{}}
\caption{\label{tab:power} Power (percent \(\pm\) MCSE) at
\(c_{bm} = 0.45\), 1000 replicates per cell.}\tabularnewline
\toprule\noalign{}
\begin{minipage}[b]{\linewidth}\raggedright
Procedure
\end{minipage} & \begin{minipage}[b]{\linewidth}\raggedright
\(N = 25\)
\end{minipage} & \begin{minipage}[b]{\linewidth}\raggedright
\(N = 35\)
\end{minipage} & \begin{minipage}[b]{\linewidth}\raggedright
\(N = 70\)
\end{minipage} & \begin{minipage}[b]{\linewidth}\raggedright
\(N = 100\)
\end{minipage} \\
\midrule\noalign{}
\endfirsthead
\toprule\noalign{}
\begin{minipage}[b]{\linewidth}\raggedright
Procedure
\end{minipage} & \begin{minipage}[b]{\linewidth}\raggedright
\(N = 25\)
\end{minipage} & \begin{minipage}[b]{\linewidth}\raggedright
\(N = 35\)
\end{minipage} & \begin{minipage}[b]{\linewidth}\raggedright
\(N = 70\)
\end{minipage} & \begin{minipage}[b]{\linewidth}\raggedright
\(N = 100\)
\end{minipage} \\
\midrule\noalign{}
\endhead
\bottomrule\noalign{}
\endlastfoot
Strict RM-ANOVA (dich.) & 35.8 \(\pm\) 1.52 & 46.8 \(\pm\) 1.58 & 75.3 \(\pm\) 1.36 & 91.3 \(\pm\) 0.89 \\
Contrast (continuous bm) & 53.8 \(\pm\) 1.58 & 67.7 \(\pm\) 1.48 & 91.5 \(\pm\) 0.88 & 98.8 \(\pm\) 0.34 \\
Linear-mixed (\texttt{corCAR1}) & 56.5 \(\pm\) 1.57 & 71.2 \(\pm\) 1.43 & 94.5 \(\pm\) 0.72 & 99.2 \(\pm\) 0.28 \\
GEE (naive sandwich) & 67.2 \(\pm\) 1.48 & 75.9 \(\pm\) 1.35 & 94.5 \(\pm\) 0.72 & 99.2 \(\pm\) 0.28 \\
GEE (Mancl-DeRouen) & 57.3 \(\pm\) 1.56 & 69.1 \(\pm\) 1.46 & 92.7 \(\pm\) 0.82 & 99.1 \(\pm\) 0.30 \\
\end{longtable}

Four substantive findings emerge from the tables. First, and most
consequentially for how the comparison should be read, the
power gap between strict RM-ANOVA and the linear-mixed model is
predominantly a cost of dichotomizing the biomarker rather than a
property of the test procedure. At \(N = 25\) the strict RM-ANOVA
power is \(0.358\); the continuous-biomarker contrast, which differs
only in retaining the continuous biomarker, reaches \(0.538\), so
dichotomization alone accounts for \(0.180\) of lost power. The
further step to the full linear-mixed model raises power only to
\(0.565\), an increment of \(0.027\) attributable to the test
procedure proper. The decomposition is stable across the grid:
the dichotomization component is \(0.209\) at \(N = 35\)
(\(0.468 \to 0.677\)), \(0.162\) at \(N = 70\) (\(0.753 \to 0.915\)),
and \(0.075\) at \(N = 100\) (\(0.913 \to 0.988\)), while the
contrast-to-linear-mixed increment never exceeds \(0.04\). The
earlier framing, which compared strict RM-ANOVA against the mixed
model directly and read the difference as a test-procedure
effect, therefore overstated the role of the test procedure by
roughly a factor of five; the honest statement is that a
classical, sphericity-free, AR(1)-free within-subject contrast on
the continuous biomarker captures most of the available power,
and the mixed model adds a small further increment.

Second, GEE with the naive sandwich is anti-conservative across
the entire sample-size range examined. Its null rejection rate is
\(0.097\) at \(N = 25\) and declines only slowly with \(N\) to
\(0.086\), \(0.079\), and \(0.074\) at \(N = 35, 70, 100\) -- still
roughly \(1.5\times\) nominal at \(N = 100\). The inflation is
therefore not a small-cluster curiosity that disappears at
trial-relevant sizes; it persists. The Mancl-DeRouen 2001
small-sample sandwich correction brings the rate within or close
to MCSE of nominal at every \(N\) (\(0.064, 0.053, 0.063, 0.059\)),
at a power cost of about ten percentage points at \(N = 25\)
(\(0.672 \to 0.573\)) and seven at \(N = 35\) (\(0.759 \to 0.691\)).
The corrected procedure's apparent power advantage over the
linear-mixed model at small \(N\) (Table \ref{tab:power}) is the
residue of its remaining slight over-rejection, not a genuine
efficiency gain.

Third, the secondary estimands (Table \ref{tab:secondary})
locate the procedures' behavior in the standard-error
calibration that the rejection rates only summarize. At
\(c_{bm} = 0.45\), the continuous-biomarker contrast and the
Mancl-DeRouen GEE are well-calibrated: their nominal-95\% CI
coverage is \(93\)--\(95\%\) and their mean model-based SE is within
a few percent of the empirical SE. Naive GEE under-covers at
small \(N\) (\(90.0\%\) at \(N = 25\), \(91.5\%\) at \(N = 35\)) because
its model-based SE is \(9\) to \(12\%\) smaller than the empirical
SE, the calibration counterpart of its Type I inflation. The
linear-mixed model over-covers (\(97.0\) to \(97.7\%\)) because its
model-based SE exceeds the empirical SE by \(11\) to \(16\%\); this
is the same conservatism visible in its sub-nominal Type I error.

\begin{longtable}[]{@{}
  >{\raggedright\arraybackslash}p{(\linewidth - 10\tabcolsep) * \real{0.2222}}
  >{\raggedright\arraybackslash}p{(\linewidth - 10\tabcolsep) * \real{0.1453}}
  >{\raggedright\arraybackslash}p{(\linewidth - 10\tabcolsep) * \real{0.1538}}
  >{\raggedright\arraybackslash}p{(\linewidth - 10\tabcolsep) * \real{0.1453}}
  >{\raggedright\arraybackslash}p{(\linewidth - 10\tabcolsep) * \real{0.1538}}
  >{\raggedright\arraybackslash}p{(\linewidth - 10\tabcolsep) * \real{0.1795}}@{}}
\caption{\label{tab:secondary} Secondary estimands at \(c_{bm} = 0.45\)
for the four coefficient-bearing arms (the interaction estimand
is \(\beta_{bm:D} = -2.21\), the Monte-Carlo reference value).
Coverage is of the nominal-95\% Wald interval; SE ratio is mean
model-based SE divided by empirical SE.}\tabularnewline
\toprule\noalign{}
\begin{minipage}[b]{\linewidth}\raggedright
Arm
\end{minipage} & \begin{minipage}[b]{\linewidth}\raggedright
Coverage \(N=25\)
\end{minipage} & \begin{minipage}[b]{\linewidth}\raggedright
Coverage \(N=100\)
\end{minipage} & \begin{minipage}[b]{\linewidth}\raggedright
SE ratio \(N=25\)
\end{minipage} & \begin{minipage}[b]{\linewidth}\raggedright
SE ratio \(N=100\)
\end{minipage} & \begin{minipage}[b]{\linewidth}\raggedright
Max \(|\)rel. bias\(|\)
\end{minipage} \\
\midrule\noalign{}
\endfirsthead
\toprule\noalign{}
\begin{minipage}[b]{\linewidth}\raggedright
Arm
\end{minipage} & \begin{minipage}[b]{\linewidth}\raggedright
Coverage \(N=25\)
\end{minipage} & \begin{minipage}[b]{\linewidth}\raggedright
Coverage \(N=100\)
\end{minipage} & \begin{minipage}[b]{\linewidth}\raggedright
SE ratio \(N=25\)
\end{minipage} & \begin{minipage}[b]{\linewidth}\raggedright
SE ratio \(N=100\)
\end{minipage} & \begin{minipage}[b]{\linewidth}\raggedright
Max \(|\)rel. bias\(|\)
\end{minipage} \\
\midrule\noalign{}
\endhead
\bottomrule\noalign{}
\endlastfoot
Contrast (continuous bm) & 94.9 & 94.8 & 0.99 & 0.98 & 3.0\% \\
Linear-mixed (\texttt{corCAR1}) & 97.7 & 97.0 & 1.16 & 1.11 & 2.6\% \\
GEE (naive sandwich) & 90.0 & 93.3 & 0.91 & 0.98 & 3.8\% \\
GEE (Mancl-DeRouen) & 94.3 & 94.0 & 1.07 & 1.02 & 3.8\% \\
\end{longtable}

Fourth, the comparison is not contaminated by point-estimate
bias. The relative bias of \(\hat{\beta}_{bm:D}\) is below \(4\%\) of
the estimand \(|\beta_{bm:D}| = 2.21\) for the contrast, the
linear-mixed model, and both GEE variants in every cell, so the
differences in power and coverage reflect variance estimation and
information use rather than systematic mis-estimation. The
linear-mixed procedure's Type I error sits below nominal
throughout (\(0.021\) to \(0.034\)), the conservatism already
visible in its over-coverage and its model/empirical SE ratio
above one; it is consistent with \texttt{corCAR1} degrees-of-freedom
accounting under the short (\(T = 8\)) serial series. A
Kenward-Roger\textsuperscript{\citeproc{ref-kenward1997small}{28}} or Satterthwaite\textsuperscript{\citeproc{ref-satterthwaite1946}{29}} small-sample degrees-of-freedom correction
would address the conservatism and recover the small power
increment it forgoes; this extension is recorded as a follow-up
item.

\section{Design-and-test recommendation}\label{design-and-test-recommendation}

The recommendation table below maps test procedure choice to
typical chronic-condition N-of-1 trial scenarios at the
prazosin/PTSD reference design (Hendrickson hybrid OL+BDC,
biomarker-treatment interaction estimand). The cycle count and
period length are held at the reference values defined in \S3
because the joint optimization over those axes is deferred to the
follow-up paper. We present the recommendation as a table of
preferred (procedure, scenario) combinations rather than as a
single `best' procedure, because the preferred choice depends on
features of the application that the analyst must determine. The
overriding recommendation, however, is design-agnostic: retain
the biomarker as a continuous predictor, since dichotomization
forfeits the large majority of the available power
(\S5, Table \ref{tab:power}).

{\def\LTcaptype{none} % do not increment counter
\begin{longtable}[]{@{}
  >{\raggedright\arraybackslash}p{(\linewidth - 4\tabcolsep) * \real{0.3032}}
  >{\raggedright\arraybackslash}p{(\linewidth - 4\tabcolsep) * \real{0.1742}}
  >{\raggedright\arraybackslash}p{(\linewidth - 4\tabcolsep) * \real{0.5226}}@{}}
\toprule\noalign{}
\begin{minipage}[b]{\linewidth}\raggedright
Scenario
\end{minipage} & \begin{minipage}[b]{\linewidth}\raggedright
Preferred procedure
\end{minipage} & \begin{minipage}[b]{\linewidth}\raggedright
Rationale
\end{minipage} \\
\midrule\noalign{}
\endhead
\bottomrule\noalign{}
\endlastfoot
Default (continuous biomarker, any \(N\)) & linear-mixed (\texttt{corCAR1}) & Highest power within nominal Type I; \texttt{corCAR1} matches AR(1) physiology; well-behaved across \(N\) \\
Lightweight / transparent classical analysis & continuous-biomarker contrast & Within 0.03 of the mixed model on power, nominal Type I, no covariance modeling required \\
\(N \le 35\), concern about LME conservatism & linear-mixed (\texttt{corCAR1}) with Kenward-Roger & LME sub-nominal and over-covers; KR addresses this without inflating Type I \\
Robust-SE-required setting (e.g.~unmodeled clustering) & GEE Mancl-DeRouen & Restores nominal Type I and 93--95\% coverage; modest power cost \\
Reviewer-mandated `classical' analysis & strict RM-ANOVA & Dichotomized sensitivity comparator; report alongside a continuous-biomarker analysis with the power-loss decomposition \\
Naive GEE (no small-sample correction) & not recommended & Type I error 1.5--1.9\(\times\) nominal across all \(N\); under-covers \\
\end{longtable}
}

The dominant message is that the biomarker should be analyzed as
a continuous predictor: dichotomizing it for a strict RM-ANOVA
forfeits the large majority of the power gap that the present
study isolates, while modeling choices among the
continuous-biomarker procedures move power only at the margin.
Within that class, the linear-mixed analysis with continuous-time
AR(1) residual correlation is the preferred default for the
biomarker-treatment interaction in aggregated N-of-1 trials at
the parameter scale evaluated here, with the dichotomization-free
within-subject contrast a close and transparent alternative that
requires no covariance modeling. Strict RM-ANOVA remains useful
as a dichotomized sensitivity comparator, as the \S2 closed-form
derivation makes explicit, but it should be reported alongside a
continuous-biomarker analysis and the power-loss decomposition
rather than on its own. Naive GEE should not be used at the
sample sizes examined without a small-sample correction; the
Mancl-DeRouen correction restores nominal Type I error and
\(93\)--\(95\%\) coverage at a moderate power cost.

\section{Discussion}\label{discussion}

\subsection{What the Monte Carlo study settles}\label{what-the-monte-carlo-study-settles}

We begin the discussion by summarizing what the Monte Carlo study
settles at the resolution it was run. The run on the 40-cell
five-procedure factorial (procedure \(\in \{\)strict RM-ANOVA,
continuous-biomarker contrast, linear-mixed with \texttt{corCAR1}, GEE
with naive sandwich, GEE with Mancl-DeRouen small-sample
correction\(\} \times c_{bm} \in \{0, 0.45\} \times N \in
\{25, 35, 70, 100\}\), with 1000 replicates per power cell and
2000 per Type I cell) was executed in 587 s with 100\%
convergence across all 60\{,\}000 fits. Per-replicate output is at
\texttt{analysis/data/quick-sim/08-test-replicates.rds}; Morris ADEMP
summary at \texttt{08-test-summary.txt}. With 1000 power replicates the
MCSE on a power estimate near \(0.5\) is approximately \(0.016\) and
with 2000 Type I replicates the MCSE on a rate near \(0.05\) is
approximately \(0.005\). The Mancl-DeRouen 2001\textsuperscript{\citeproc{ref-mancl2001covariance}{24}} correction was implemented inline by
inflating each cluster's residual via \((I - H_{ii})^{-1} r_i\) in
the sandwich meat; it produces identical point estimates and
larger standard errors than the naive sandwich, and its SE for
the interaction coefficient agrees with \texttt{geesmv::GEE.var.md} to
within approximately 2.5\% (mean inline/\texttt{geesmv} ratio 1.02 over
25 probe replicates), so the correction magnitude is externally
validated.

\section{Conclusions}\label{conclusions}

In conclusion, three points are to be emphasized. First, the
analyst's most consequential choice for the biomarker-treatment
interaction in aggregated N-of-1 trials is to keep the biomarker
continuous rather than to dichotomize it for a classical
RM-ANOVA. The continuous-biomarker contrast recovers about
four-fifths of the apparent RM-ANOVA-to-mixed-model power gap
(dichotomization costs \(0.08\) to \(0.21\) of power across the
sample sizes evaluated, while the test procedure proper adds at
most \(0.04\)); the linear-mixed analysis with continuous-time
AR(1) residual correlation has the highest power within nominal
Type I among the continuous-biomarker procedures, with the
dichotomization-free within-subject contrast a close alternative.
Naive GEE with the standard sandwich estimator inflates Type I
error to between \(1.5\times\) and \(1.9\times\) nominal across the
whole \(N \in \{25,35,70,100\}\) range, disqualifying its rejection
rate as a power summary; the Mancl-DeRouen 2001 small-sample
sandwich correction restores nominal Type I control and
\(93\)--\(95\%\) coverage at a moderate (about 10-percentage-point)
power cost, making it the operational GEE comparator rather than
the naive sandwich.

Second, the closed-form strict-RM-ANOVA test of the
biomarker-treatment interaction (\S2) reduces to a two-sample
\(t\)-test on the within-subject treatment differences under
sphericity and balance. The reduction makes the strict-RM-ANOVA
result analytically transparent and connects the simulation
findings to the classical split-plot literature. It also pins
down what does and does not limit the test here. The strict
RM-ANOVA arm we evaluate pre-averages each phase to a single
value, so its within-subject covariance is \(2 \times 2\) and
sphericity holds by construction; consistent with this, its Type
I error is at nominal (Table \ref{tab:type1}), and its power
deficit is wholly attributable to the median dichotomization of
the biomarker, as the continuous-biomarker contrast demonstrates.
The sphericity penalty that motivates much of \S2 applies instead
to the alternative classical analysis that treats the full
\(T = 8\) series as repeated within-subject levels: there the AR(1)
residual correlation violates sphericity and a
Greenhouse-Geisser correction recovers nominal Type I only at a
degrees-of-freedom cost, which the linear-mixed analysis with
\texttt{corCAR1} avoids by modeling the AR(1) structure directly. Both
considerations point the same way -- away from a dichotomized
classical analysis -- but through different mechanisms, and we
are careful to keep them distinct.

Third, the cycle-by-period design grid that constitutes the
second axis of the joint design-by-test sensitivity question
framed in \S1 is deferred to a follow-up paper. The
infrastructure for the sweep is in place at
\texttt{analysis/scripts/test-procedure-design-sensitivity/}, and the
design specification is committed at
\texttt{docs/27-cycle-period-design-sweep-plan.md}. The follow-up paper
will deliver the joint (cycle count, period length, test
procedure) recommendation table that the present paper's \S6
narrows to the test-procedure-only axis.

\section{Conflict of interest}\label{conflict-of-interest}

The authors declare no conflicts of interest.

\section{Funding}\label{funding}

This work received no specific funding.

\section{Data availability statement}\label{data-availability-statement}

The data and code that support the findings of this study are
openly available in the pmsimstats-ng project repository.
Per-replicate Monte Carlo outputs and ADEMP-compliant cell-level
summaries are versioned under \texttt{analysis/data/}; simulation
drivers are at \texttt{analysis/scripts/}. Exact paths relevant to this
manuscript are listed in the Reproducibility section below.

\section{Reproducibility}\label{reproducibility}

This manuscript's simulation was executed with R 4.6.x; the
project also ships a zzcollab Docker container (image hash
recorded in \texttt{Dockerfile}; package set captured in \texttt{renv.lock})
for container-based reproduction. Principal packages used by the
Study 1 pipeline are \texttt{nlme} (continuous-time linear-mixed
analysis with \texttt{corCAR1} residual correlation), \texttt{geepack} (GEE
fitting) plus an inline Mancl-DeRouen 2001 small-sample sandwich
correction, \texttt{geesmv} (independent benchmark of that correction),
\texttt{data.table}, and \texttt{parallel} (\texttt{mclapply} for the parameter-grid
sweep), with \texttt{rmarkdown} plus \texttt{bookdown} for the manuscript
build. The Study 1 simulation runs on the pmsimstats package
data-generating process under Architecture A via \texttt{generateData()}
and \texttt{lme\_analysis()}, not on the reduced
\texttt{implementations/simple} substrate.

\textbf{Random-number generator.} The driver sets
\texttt{RNGkind("L\textquotesingle{}Ecuyer-CMRG")} and \texttt{set.seed(20260507)}, then runs
the cell-by-replicate jobs through \texttt{parallel::mclapply} with
\texttt{mc.set.seed\ =\ TRUE}, so each replicate draws from an
independent, reproducible random-number stream.

\textbf{Data and code availability.} Per-replicate simulation outputs
are checkpointed at
\texttt{analysis/data/quick-sim/08-test-replicates.rds}; the extended
Morris ADEMP cell-level summary (rejection rate, MCSE,
convergence, bias, empirical and model SE, coverage) at the
companion \texttt{08-test-summary.txt}, with the Monte-Carlo reference
estimand at \texttt{08-test-betatrue.txt} and the Mancl-DeRouen
benchmark at \texttt{08-test-mancl-benchmark.txt}. The driver is
\texttt{analysis/scripts/test-procedure-design-sensitivity/01-study1-test-procedure.R}
(superseding the earlier
\texttt{analysis/scripts/quick-sim/08-test-procedure-prototype.R}). The
pre-registered ADEMP plan, and the deviations recorded in Table
\ref{tab:deviations}, are at
\texttt{analysis/scripts/test-procedure-design-sensitivity/00-ademp-pre-registration.md}.
The manuscript source, paper-specific bibliography
(\texttt{references.bib}), and SIM CSL file are at
\texttt{analysis/report/08-test-procedure-design-sensitivity/}.

\section{References}\label{references}

\protect\phantomsection\label{refs}
\begin{CSLReferences}{0}{1}
\bibitem[\citeproctext]{ref-morris2019}
\CSLLeftMargin{1. }%
\CSLRightInline{Morris TP, White IR, Crowther MJ. Using simulation studies to evaluate statistical methods. \emph{Statistics in Medicine}. 2019;38(11):2074-2102.}

\bibitem[\citeproctext]{ref-haverkamp2017sphericity}
\CSLLeftMargin{2. }%
\CSLRightInline{Haverkamp N, Beauducel A. Violation of the sphericity assumption and its effect on type-{I} error rates in repeated measures {ANOVA} and multi-level linear models. \emph{Frontiers in Psychology}. 2017;8:1841. doi:\href{https://doi.org/10.3389/fpsyg.2017.01841}{10.3389/fpsyg.2017.01841}}

\bibitem[\citeproctext]{ref-huynh1976}
\CSLLeftMargin{3. }%
\CSLRightInline{Huynh H, Feldt LS. Estimation of the {Box} correction for degrees of freedom from sample data in randomized block and split-plot designs. \emph{Journal of Educational Statistics}. 1976;1(1):69-82. doi:\href{https://doi.org/10.3102/10769986001001069}{10.3102/10769986001001069}}

\bibitem[\citeproctext]{ref-greenhouse1959}
\CSLLeftMargin{4. }%
\CSLRightInline{Greenhouse SW, Geisser S. On methods in the analysis of profile data. \emph{Psychometrika}. 1959;24(2):95-112. doi:\href{https://doi.org/10.1007/BF02289823}{10.1007/BF02289823}}

\bibitem[\citeproctext]{ref-blanca2023sphericity}
\CSLLeftMargin{5. }%
\CSLRightInline{Blanca MJ, Arnau J, Garcı́a-Castro FJ, Alarcón R, Bono R. Repeated measures {ANOVA} and adjusted tests when sphericity is violated: Which procedure is best? \emph{Frontiers in Psychology}. 2023;14:1192453. doi:\href{https://doi.org/10.3389/fpsyg.2023.1192453}{10.3389/fpsyg.2023.1192453}}

\bibitem[\citeproctext]{ref-natesanbatley2023}
\CSLLeftMargin{6. }%
\CSLRightInline{Natesan Batley P, McClure EB, Brewer B, et al. Evidence and reporting standards in {N}-of-1 medical studies: A systematic review. \emph{Translational Psychiatry}. 2023;13(1):263. doi:\href{https://doi.org/10.1038/s41398-023-02562-8}{10.1038/s41398-023-02562-8}}

\bibitem[\citeproctext]{ref-stunnenberg2022}
\CSLLeftMargin{7. }%
\CSLRightInline{Stunnenberg BC, Berends J, Griggs RC, et al. {N}-of-1 trials in neurology: A systematic review. \emph{Neurology}. 2022;98(2):e174-e185. doi:\href{https://doi.org/10.1212/WNL.0000000000012998}{10.1212/WNL.0000000000012998}}

\bibitem[\citeproctext]{ref-shaffer2018}
\CSLLeftMargin{8. }%
\CSLRightInline{Shaffer JA, Kronish IM, Falzon L, Cheung YK, Davidson KW. {N}-of-1 randomized intervention trials in health psychology: A systematic review and methodology critique. \emph{Annals of Behavioral Medicine}. 2018;52(9):731-742. doi:\href{https://doi.org/10.1093/abm/kax026}{10.1093/abm/kax026}}

\bibitem[\citeproctext]{ref-vohra2015cent}
\CSLLeftMargin{9. }%
\CSLRightInline{{Vohra S, Shamseer L, Sampson M, et al.} {CONSORT} extension for reporting {N}-of-1 trials ({CENT}) 2015 statement. \emph{Journal of Clinical Epidemiology}. 2015;76:9-17. doi:\href{https://doi.org/10.1016/j.jclinepi.2015.05.004}{10.1016/j.jclinepi.2015.05.004}}

\bibitem[\citeproctext]{ref-shamseer2015cent}
\CSLLeftMargin{10. }%
\CSLRightInline{{Shamseer L, Sampson M, Bukutu C, Schmid CH, et al.} {CENT} 2015: Explanation and elaboration. \emph{Journal of Clinical Epidemiology}. 2015;76:18-46. doi:\href{https://doi.org/10.1016/j.jclinepi.2015.05.018}{10.1016/j.jclinepi.2015.05.018}}

\bibitem[\citeproctext]{ref-kwasnicka2019}
\CSLLeftMargin{11. }%
\CSLRightInline{{Kwasnicka D, Inauen J, Nieuwenboom W, et al.} Challenges and solutions for {N}-of-1 design studies in health psychology. \emph{Health Psychology Review}. 2019;13(2):163-178. doi:\href{https://doi.org/10.1080/17437199.2018.1564627}{10.1080/17437199.2018.1564627}}

\bibitem[\citeproctext]{ref-ma2012beyond}
\CSLLeftMargin{12. }%
\CSLRightInline{Ma Y, Mazumdar M, Memtsoudis SG. Beyond repeated-measures analysis of variance: Advanced statistical methods for the analysis of longitudinal data in anesthesia research. \emph{Regional Anesthesia and Pain Medicine}. 2012;37(1):99-105. doi:\href{https://doi.org/10.1097/AAP.0b013e31823ebc74}{10.1097/AAP.0b013e31823ebc74}}

\bibitem[\citeproctext]{ref-mallinckrodt2003}
\CSLLeftMargin{13. }%
\CSLRightInline{Mallinckrodt CH, Clark SWS, Carroll RJ, Molenberghs G. Assessing response profiles from incomplete longitudinal clinical trial data under regulatory considerations. \emph{Journal of Biopharmaceutical Statistics}. 2003;13(2):179-190. doi:\href{https://doi.org/10.1081/BIP-120019265}{10.1081/BIP-120019265}}

\bibitem[\citeproctext]{ref-mallinckrodt2001dropout}
\CSLLeftMargin{14. }%
\CSLRightInline{Mallinckrodt CH, Clark WS, David SR. Accounting for dropout bias using mixed-effects models. \emph{Journal of Biopharmaceutical Statistics}. 2001;11(1-2):9-21. doi:\href{https://doi.org/10.1081/BIP-100104194}{10.1081/BIP-100104194}}

\bibitem[\citeproctext]{ref-mallinckrodt2004duloxetine}
\CSLLeftMargin{15. }%
\CSLRightInline{Mallinckrodt CH, Raskin J, Wohlreich MM, Watkin JG, Detke MJ. The efficacy of duloxetine: A comprehensive summary of results from {MMRM} and {LOCF\_ANCOVA} in eight clinical trials. \emph{BMC Psychiatry}. 2004;4:26. doi:\href{https://doi.org/10.1186/1471-244X-4-26}{10.1186/1471-244X-4-26}}

\bibitem[\citeproctext]{ref-prakash2008}
\CSLLeftMargin{16. }%
\CSLRightInline{Prakash A, Risser RC, Mallinckrodt CH. The impact of analytic method on interpretation of outcomes in longitudinal clinical trials. \emph{International Journal of Clinical Practice}. 2008;62(8):1147-1158. doi:\href{https://doi.org/10.1111/j.1742-1241.2008.01808.x}{10.1111/j.1742-1241.2008.01808.x}}

\bibitem[\citeproctext]{ref-liang1986}
\CSLLeftMargin{17. }%
\CSLRightInline{Liang KY, Zeger SL. Longitudinal data analysis using generalized linear models. \emph{Biometrika}. 1986;73(1):13-22. doi:\href{https://doi.org/10.1093/biomet/73.1.13}{10.1093/biomet/73.1.13}}

\bibitem[\citeproctext]{ref-wang2016geesmv}
\CSLLeftMargin{18. }%
\CSLRightInline{Wang M, Kong L, Li Z, Zhang L. Covariance estimators for generalized estimating equations ({GEE}) in longitudinal analysis with small samples. \emph{Statistics in Medicine}. 2016;35(10):1706-1721. doi:\href{https://doi.org/10.1002/sim.6817}{10.1002/sim.6817}}

\bibitem[\citeproctext]{ref-pan2001}
\CSLLeftMargin{19. }%
\CSLRightInline{Pan W. Sample size and power calculations with correlated binary data. \emph{Controlled Clinical Trials}. 2001;22(3):211-227. doi:\href{https://doi.org/10.1016/s0197-2456(01)00131-3}{10.1016/s0197-2456(01)00131-3}}

\bibitem[\citeproctext]{ref-tong2024stepwedge}
\CSLLeftMargin{20. }%
\CSLRightInline{{Tong G, Nevins P, Ryan M, others, Taljaard M}. A review of current practice in the design and analysis of extremely small stepped-wedge cluster randomized trials. \emph{Clinical Trials}. 2024;22(1):45-56. doi:\href{https://doi.org/10.1177/17407745241276137}{10.1177/17407745241276137}}

\bibitem[\citeproctext]{ref-defelippe2023}
\CSLLeftMargin{21. }%
\CSLRightInline{{Defelippe VM, Thiel GJMW van, Otte WM, Schutgens REG, Stunnenberg B, et al.} Toward responsible clinical {n}-of-1 strategies for rare diseases. \emph{Drug Discovery Today}. 2023;28(10):103688. doi:\href{https://doi.org/10.1016/j.drudis.2023.103688}{10.1016/j.drudis.2023.103688}}

\bibitem[\citeproctext]{ref-senn1994crossover}
\CSLLeftMargin{22. }%
\CSLRightInline{Senn S. The {AB/BA} crossover: Past, present and future? \emph{Statistical Methods in Medical Research}. 1994;3(4):303-324. doi:\href{https://doi.org/10.1177/096228029400300402}{10.1177/096228029400300402}}

\bibitem[\citeproctext]{ref-senn2002crossover}
\CSLLeftMargin{23. }%
\CSLRightInline{Senn S. \emph{Cross-over Trials in Clinical Research}. 2nd ed. Wiley; 2002.}

\bibitem[\citeproctext]{ref-mancl2001covariance}
\CSLLeftMargin{24. }%
\CSLRightInline{Mancl LA, DeRouen TA. A covariance estimator for {GEE} with improved small-sample properties. \emph{Biometrics}. 2001;57(1):126-134.}

\bibitem[\citeproctext]{ref-maxwelldelaney2018}
\CSLLeftMargin{25. }%
\CSLRightInline{Maxwell SE, Delaney HD, Kelley K. \emph{Designing Experiments and Analyzing Data: A Model Comparison Perspective}. 3rd ed. Routledge; 2018.}

\bibitem[\citeproctext]{ref-mallinckrodt2008recommendations}
\CSLLeftMargin{26. }%
\CSLRightInline{Mallinckrodt CH, Lane PW, Schnell D, Peng Y, Mancuso JP. Recommendations for the primary analysis of continuous endpoints in longitudinal clinical trials. \emph{Drug Information Journal}. 2008;42(4):303-319. doi:\href{https://doi.org/10.1177/009286150804200402}{10.1177/009286150804200402}}

\bibitem[\citeproctext]{ref-hendrickson2020}
\CSLLeftMargin{27. }%
\CSLRightInline{Hendrickson RC, Thomas RG, Schork NJ, Raskind MA. Optimizing aggregated {N}-of-1 trial designs for predictive biomarker validation: Statistical methods and practical considerations. \emph{Frontiers in Digital Health}. 2020;2:13. doi:\href{https://doi.org/10.3389/fdgth.2020.00013}{10.3389/fdgth.2020.00013}}

\bibitem[\citeproctext]{ref-kenward1997small}
\CSLLeftMargin{28. }%
\CSLRightInline{Kenward MG, Roger JH. Small sample inference for fixed effects from restricted maximum likelihood. \emph{Biometrics}. 1997;53(3):983-997. doi:\href{https://doi.org/10.2307/2533558}{10.2307/2533558}}

\bibitem[\citeproctext]{ref-satterthwaite1946}
\CSLLeftMargin{29. }%
\CSLRightInline{Satterthwaite FE. An approximate distribution of estimates of variance components. \emph{Biometrics Bulletin}. 1946;2(6):110-114. doi:\href{https://doi.org/10.2307/3002019}{10.2307/3002019}}

\end{CSLReferences}

\vfill

\begin{center}\rule{0.5\linewidth}{0.5pt}\end{center}

\emph{Rendered on 2026-07-30 at 14:14 PDT.}\\
\emph{Source: \texttt{\textasciitilde{}/prj/res/36-pmsimstats-ng/pmsimstats-ng/analysis/report/08-test-procedure-design-sensitivity/report-slim.Rmd}}

\end{document}
